% SEGMENT OLP-0630-S01
% Part: history
% Chapter: biographies
% Section: alfred-tarski

\documentclass[../../../include/open-logic-section]{subfiles}

\begin{document}

\olfileid{his}{bio}{tar}

\olsection{ஆல்ஃபிரட் தார்ஸ்கி}

\olphoto{tarski-alfred}{ஆல்ஃபிரட் தார்ஸ்கி}

ஆல்ஃபிரட் தார்ஸ்கி 1901 ஜனவரி 14 அன்று போலந்தின் வார்சாவில் பிறந்தார்;
அந்நகரம் அப்போது ரஷ்யப் பேரரசின் பகுதியாக இருந்தது. ``நெப்போலியன்
போன்றவர்'' என்று வருணிக்கப்பட்ட தார்ஸ்கி ஆரவாரமும் பேச்சும் தீவிரமும்
மிக்கவர். இந்த ஆற்றல் அவரது விரிவுரைகளிலும் வெளிப்பட்டது. ஒருமுறை
விரிவுரையின்போது சிகரெட்டை எறிந்தபோது குப்பைக் கூடைக்குத் தீ
பிடித்தது; பின்னர் அந்தக் கட்டடத்தில் விரிவுரையாற்ற அவருக்குத்
தடை விதிக்கப்பட்டது.

% SEGMENT OLP-0630-S02
இளமையிலிருந்தே தார்ஸ்கி அறிவை நாடியவர். பிற்காலத்தில், தருக்கவியலில்
மட்டுமே தமக்கு B தர மதிப்பெண் கிடைத்ததால் அதைப் படித்ததாக மாணவர்களிடம்
கூறுவார். ஆனால் பள்ளிப் பதிவுகள் தருக்கவியல் உட்பட எல்லாப்
பாடங்களிலும் அவர் A தர மதிப்பெண் பெற்றதைக் காட்டுகின்றன. 1918 முதல்
1924 வரை வார்சா பல்கலைக்கழகத்தில் படித்தார். முதலில் உயிரியல்
படிக்க எண்ணினார்; ஆனால் அந்தப் பல்கலைக்கழகம் வார்சா தருக்கவியல்,
மெய்யியல் பள்ளியின் மையமாக இருந்ததால் கணிதம், மெய்யியல்,
தருக்கவியல் ஆகியவற்றில் ஆர்வம் கொண்டார். ஸ்டானிஸ்லாவ் லெஷ்னெவ்ஸ்கியின்
(Stanis\l{}aw Le\'{s}niewski) மேற்பார்வையில் 1924-இல் முனைவர்
பட்டம் பெற்றார்.

% SEGMENT OLP-0630-S03
1939-இல் அமெரிக்க ஐக்கிய நாடுகளுக்குக் குடிபெயர்வதற்கு முன்,
வார்சாவில் மேல்நிலைப் பள்ளி ஆசிரியராகப் பணியாற்றிய காலத்திலேயே
தார்ஸ்கி தமது மிக முக்கியமான சில ஆய்வுகளைச் செய்தார். தருக்கரீதியான
பின்விளைவு, தருக்கரீதியான மெய்மை பற்றிய எழுத்துகள் அக்காலத்தவை.
1939-இல் விரிவுரைச் சுற்றுப்பயணத்திற்காக அமெரிக்காவுக்கு வந்திருந்தார்.
அவர் அங்கிருந்தபோது ஜெர்மனி போலந்தின்மீது படையெடுத்தது; தாம் யூதர்
என்பதால் அவர் திரும்ப முடியவில்லை. அவருடைய மனைவியும் குழந்தைகளும்
போர் முடியும்வரை போலந்திலேயே இருந்தனர்; பின்னர் அவர்களும்
அமெரிக்காவுக்குக் குடிபெயர முடிந்தது. தார்ஸ்கி ஹார்வர்ட்,
நியூ யார்க் நகரக் கல்லூரி, பிரின்ஸ்டனின் மேம்பட்ட ஆய்வுகள் நிறுவனம்
ஆகியவற்றிலும், இறுதியாக கலிபோர்னியா பல்கலைக்கழகத்தின் பெர்க்லி
வளாகத்திலும் கற்பித்தார். அங்கே தருக்கவியலும் அறிவியலின்
முறைமையியலும் குறித்த பல்துறைத் திட்டத்தை நிறுவினார். 1983
அக்டோபர் 26 அன்று 82 வயதில் இறந்தார்.

% SEGMENT OLP-0630-S04
\begin{reading}
தார்ஸ்கியின் வாழ்க்கை பற்றி மேலும் அறிய \emph{Alfred Tarski: Life
  and Logic} என்ற வாழ்க்கை வரலாற்றைப் பார்க்கவும் \citep{Feferman2004}.
தருக்கரீதியான பின்விளைவு, மெய்மை பற்றிய அவருடைய முக்கியமான
எழுத்துகள் ஆங்கிலத்தில் \citep{Tarski1983} இல் கிடைக்கின்றன.
மூல எழுத்துகள் அனைத்தும் நான்கு தொகுதிகளாகச் சேர்க்கப்பட்டுள்ளன
\citep{Tarski1981}.
\end{reading}

\end{document}
